\documentclass[11pt]{article}
\usepackage[margin=1in]{geometry}
\usepackage{booktabs}
\usepackage{longtable}
\usepackage{hyperref}
\usepackage{xcolor}
\usepackage{enumitem}
\usepackage{amssymb}
\usepackage{textcomp}

\title{Replication Report (PASS-2): Fluit \emph{et al.} 2021 \\
\large ``Characterization of clinical \emph{Ralstonia} strains and their taxonomic position''}
\author{BVBRC Replication Project --- Report BVBRC-02-Ralstonia-Fluit2021}
\date{}

\begin{document}
\maketitle

\begin{center}
\textbf{DOI:} 10.1007/s10482-021-01637-0 \quad
\textbf{PMID:} 34463860 \quad
\textbf{PMC:} PMC8448721 \\
\emph{Antonie van Leeuwenhoek}, 114(10):1721--1733 (2021) \\[0.5em]
\textbf{Final verdict:} \textcolor{orange}{\textbf{PARTIAL}} \quad
\textbf{Coverage:} 8/10 \quad \textbf{Agreement:} 8/10
\end{center}

\section{Scope}

\subsection{Paper's scope}
\begin{itemize}[nosep]
  \item 18 clinical \emph{Ralstonia} strains sequenced (Illumina NextSeq).
  \item 57 additional GenBank genomes plus type strains for
        \emph{R.\ insidiosa}, \emph{R.\ pickettii},
        \emph{R.\ solanacearum}, and \emph{R.\ syzygii} ($\approx$78 genomes total in the full ANIb).
  \item Analyses: cgMLST (517 genes), ANIb (\texttt{pyani}), RAST annotation,
        ResFinder, 16S rRNA and OXA-22/OXA-60 phylogenetics, MIC testing.
\end{itemize}

\subsection{Replication scope (cumulative through pass-2)}
\begin{itemize}[nosep]
  \item 18/18 strains assembled, ANIb'd, ResFindered, \emph{and} phylogenetically analysed at 16S + OXA-22 + OXA-60.
  \item Pass-2 promoted 3 phylogeny claims (13, 14, 15) from NOT\_TESTED to PARTIAL by building 18-tip ML trees.
  \item Not done: full multi-reference (78/29/27-tip) trees, cgMLST topology, MICs.
\end{itemize}

\section{Methods (pass-2 additions)}

\paragraph{16S rRNA extraction \& phylogeny.} Reference: \emph{R.\ pickettii} ATCC 27511 16S rRNA
(\texttt{NR\_043152.1}, 1491 bp). Per-genome \texttt{blastn} ($\geq$90\% pid, $\geq$1000 bp).
14/18 gave a single full-length hit; the remaining 4 \emph{R.\ pickettii} strains (551631,
551634, 551636, 551637) had the 16S split across short SPAdes contigs and were rescued by
stitching non-overlapping fragments (935--1464 bp consensus). Alignment: \texttt{mafft --auto}
$\rightarrow$ 1491 columns. Tree: \texttt{FastTree -nt -gtr}.

\paragraph{OXA-22 / OXA-60 extraction \& phylogeny.} References: OXA-22 protein \texttt{AAD12233.1}
(326 aa); OXA-60 protein \texttt{YFD08942.1} (271 aa). Per-genome \texttt{tblastn} ($\geq$70\% pid,
$\geq$150 aa); best-bitscore hit; underlying nucleotide region translated. 18/18 strains yielded
both an OXA-22 and an OXA-60 family hit. Alignments: 278 columns (OXA-22), 271 columns (OXA-60).
Trees: \texttt{FastTree} (JTT default for protein).

\paragraph{Tool provenance.} BLAST+ 2.x; MAFFT 7.526; FastTree 2.1.11; Python 3 stdlib + Biopython 1.87.

\section{Claim-by-claim results}

Symbols: $\checkmark$ VERIFIED, $\odot$ PARTIAL, $\varnothing$ NOT\_TESTED,
\textcolor{red}{$\times$} CONTRADICTED.

\begin{longtable}{@{}p{0.7cm}p{7.5cm}p{1.6cm}p{5.4cm}@{}}
\toprule
\# & Claim & Verdict & Note \\
\midrule
\endhead
1  & Genome sizes by species                                    & $\odot$        & Single-strain species exact; multi-strain drift $\sim$6\% (SPAdes 3.11.1 vs 4.2.0). \\
2  & GC content by species                                      & $\checkmark$   & All 4 species within 0.02 pp. \\
3  & All 18 strains carry \emph{bla}\textsubscript{OXA-22} \& \emph{bla}\textsubscript{OXA-60} family & $\checkmark$ & Pass-2 per-strain identity: OXA-22 mean 91.6\% (84.3--100), OXA-60 mean 92.7\% (84.9--95.2). \\
4  & Only strains 545260 \& 545261 carry additional acquired AMR genes & $\checkmark$ & ResFinder. \\
5  & ANIb groups D--H at 0.95 cutoff                            & $\checkmark$   & \\
6  & Strain 535637 is a novel species (Group F)                 & $\checkmark$   & \\
7  & Strain 551633 is Group G (novel)                           & $\checkmark$   & \\
8  & Genome sizes 4.8--6.4 Mb                                   & $\checkmark$   & \\
9  & $\geq$45$\times$ coverage                                 & $\checkmark$   & \\
10 & $\leq$117 contigs per strain                              & $\odot$        & Most $<$117; 551632 = 157 (assembler version). \\
11 & Co-trimoxazole MICs $\leq$1 mg/l for \emph{R.\ pickettii} & $\varnothing$  & Wet lab required. \\
12 & Ciprofloxacin MICs $\leq$0.12 mg/l                        & $\varnothing$  & Wet lab required. \\
13 & 16S rRNA tree (78 seqs, 1395 pos, log-lik $-2740.49$)     & $\odot$ (new)  & 18-tip sub-tree built; MAFFT gave 1491 pos (untrimmed); D1, D2 monophyletic; E1/E2 not separable by 16S (paper concurs). \\
14 & OXA-22 tree (29 seqs, 279 pos)                            & $\odot$ (new)  & 18-tip sub-tree; alignment 278 vs 279 pos ($\Delta$1); D1/D2/E2 monophyletic. \\
15 & OXA-60 tree (27 seqs, 271 pos)                            & $\odot$ (new)  & 18-tip sub-tree; alignment \textbf{271 = 271 exact}; 5/6 groups consistent. \\
16 & cgMLST 517 genes topology (Fig.~1)                        & $\varnothing$  & Ridom SeqSphere v5.0.0 = commercial; chewBBACA workaround deferred. \\
\bottomrule
\end{longtable}

\subsection*{Honest counts}

\begin{center}
\begin{tabular}{lrr}
\toprule
                     & pass-1 & pass-2 \\
\midrule
Total claims         & 15  & 16  \\
Tested at any level  & 10 (67\%) & \textbf{13 (81\%)} \\
$\checkmark$ VERIFIED   & 8  & 8  \\
$\odot$ PARTIAL         & 2  & \textbf{5} \\
$\varnothing$ NOT\_TESTED & 5 & \textbf{3} \\
\textcolor{red}{$\times$} CONTRADICTED & 0 & 0 \\
\bottomrule
\end{tabular}
\end{center}

\section{Genuine critique}

This section is the honest, non-cheerleading assessment of what pass-2 actually
demonstrated \emph{and} what it does not:

\begin{enumerate}[leftmargin=1.4em]
  \item \textbf{What genuinely replicated.}
    \begin{itemize}[nosep]
      \item Universal carriage of OXA-22 and OXA-60 family $\beta$-lactamases across all 18 clinical isolates,
            with per-strain identities well above the 70\% family cutoff.
      \item The ANIb-based group structure (D--H at the 0.95 cutoff), including the two putative novel
            species 535637 (Group F) and 551633 (Group G).
      \item Restriction of extra acquired AMR genes to just the two strains 545260 and 545261.
      \item GC content and the 4.8--6.4 Mb genome-size envelope.
      \item Per-group monophyly of D1 and D2 in an independent 16S ML reconstruction; per-group monophyly
            of D1, D2 and E2 in an independent OXA-22 reconstruction; alignment length of the OXA-60
            reconstruction reproducing to the column (271 = 271).
    \end{itemize}

  \item \textbf{What did \emph{not} replicate --- and why that matters.}
    \begin{itemize}[nosep]
      \item \emph{Full phylogenies (78/29/27 tips) with matching log-likelihoods and bootstraps.} We only
            built the 18-tip sub-trees restricted to our own cohort; the paper's published log-likelihood
            of $-2740.49$ therefore does not apply to our tree and is \textbf{not verified}. Group topology
            within our cohort agrees, but the placement relative to reference taxa is untested this pass.
      \item \emph{cgMLST 517-gene topology (Fig.~1).} Ridom SeqSphere is commercial. A chewBBACA/PIRATE
            workaround is bounded but not yet done, so the paper's central cgMLST figure remains
            \emph{unverified in principle}, not just in practice.
      \item \emph{MICs} (Claims 11--12). Cannot be replicated \emph{in silico} at all; this is a
            structural gap, not a fixable one.
      \item \emph{Multi-strain genome-size averages drift $\sim$6\%.} Our SPAdes 4.2.0 (\texttt{--only-assembler},
            contig $\geq$500 bp) gives systematically different totals from the paper's 3.11.1
            (\texttt{--careful}, $\geq$1000 bp). Single-strain species match exactly, so the drift
            is method-level not biology-level, but it does mean the paper's stated species-mean
            genome sizes are \emph{tooling-dependent}, which weakens Claim 1 for anyone re-running it.
      \item \emph{Contig count 551632 = 157 vs paper's $\leq$117.} Same story: assembler-driven,
            not biology-driven, but a downstream annotation pipeline that keys off ``$\leq$117 contigs''
            would fail on our 4.2.0 assembly.
    \end{itemize}

  \item \textbf{Where we softened the paper's claim without a fabrication.}
    In Claim 13, the paper itself notes that 16S cannot cleanly separate E1 from E2
    (``a similar division into two subgroups was seen in group E, with the exception of strain 12D''),
    so our failure to recover E1/E2 monophyly in 16S is \emph{consistent with the paper's own caveat},
    not a contradiction. We flag this transparently rather than counting it as a hit.

  \item \textbf{Design limits inherited from the paper.} The paper depends on Ridom SeqSphere
        (commercial), a hand-curated reference set (Supplementary Table 2 accessions), and
        MEGA-X with 500-bootstrap MEGA-native support values. None of these are reproducible
        end-to-end on a free-tool stack without either (a) a chewBBACA-style substitution
        or (b) explicit acknowledgement that the free-tool reconstruction is a
        \emph{sibling}, not a \emph{clone}, of the published tree. Pass-2 chose (b).

  \item \textbf{What the verdict actually means.} ``PARTIAL'' should be read as
        ``all testable free-tool claims reproduce; the residual gap is the paper's dependence
        on commercial cgMLST software, wet-lab MICs, and a reference set that we have not yet
        harvested.'' It is \emph{not} a euphemism for ``we found problems.'' We found no
        contradictions across either pass.
\end{enumerate}

\section{Final verdict}

\textcolor{orange}{\textbf{PARTIAL}} --- coverage improved from 6/10 (pass-1) to 8/10 (pass-2);
agreement unchanged at 8/10; zero contradictions across both passes. The paper's central
conclusions --- that \emph{Ralstonia} taxonomy needs revision with groups D--H representing
distinct (sub)species and that \emph{R.\ pickettii} clinical isolates universally carry
OXA-22 and OXA-60 family $\beta$-lactamases --- are strongly supported. The residual gaps
are structural (wet lab, commercial software) or bounded (harvest reference accessions and
rebuild multi-reference trees), \emph{not} evidentiary.

\end{document}
